% Figure — Contre-réaction courant-série : amplificateur A et réseau de contre-réaction β_F (EANA3, chapitre 2).
% D'après les notes externes [notes-2023-c04-06, p. 1], [tableau-2020-c03, p. 6] et [tableau-2020-c04, p. 1], redessinée.
% Rendu SVG :
%   latex -output-directory=<dossier temporaire> 02-cr-courant-serie.tex
%   dvisvgm --no-fonts --exact-bbox -o 02-cr-courant-serie.svg <dossier temporaire>/02-cr-courant-serie.dvi
\documentclass[border=4pt]{standalone}
\usepackage[T1]{fontenc}
\usepackage[utf8]{inputenc}
\usepackage{amsmath}
\usepackage{sansmath}
\renewcommand{\familydefault}{\sfdefault}
\usepackage[european, straightvoltages, americancurrents]{circuitikz}
\begin{document}
\sansmath
\begin{circuitikz}[font=\small, >=latex, fleche/.style={->, thick}]
  \draw[thick, fill=black!3] (1.9,1.6) rectangle (7.4,4.3); \node[font=\bfseries] at (7.1,4.05) {$A$};
  \draw[thick, fill=black!3] (1.9,-1.7) rectangle (7.4,0.6); \node at (4.8,-1.45) {$\beta_F$};
  \draw (1.9,3.5) node[circ]{} -- (2.8,3.5) to[generic, l=$Z_e$] (2.8,2.0) -- (1.9,2.0) node[circ]{};
  \draw[fleche] (2.25,2.15) -- node[left]{$v_e$} (2.25,3.35);
  \draw (4.6,3.5) to[american controlled current source, l=$A v_e$] (4.6,2.0);
  \draw (6.5,3.5) to[generic, l=$Z_s$] (6.5,2.0);
  \draw (4.6,3.5) -- (7.4,3.5) node[circ]{};  \draw (4.6,2.0) -- (7.4,2.0) node[circ]{};
  \draw (1.9,0.2) node[circ]{} -- (3.0,0.2) to[american controlled voltage source, l=$\beta_F\, i_o$] (3.0,-1.2) -- (1.9,-1.2) node[circ]{};
  \draw[fleche] (2.3,-1.05) -- node[left]{$v_{fb}$} (2.3,0.05);
  \draw (7.4,0.2) node[circ]{} -- (6.4,0.2) -- (6.4,-1.2) -- (7.4,-1.2) node[circ]{};
  \draw[fleche] (6.25,-0.05) -- node[left]{$i_o$} (6.25,-0.95);
  \draw (0,3.5) node[ocirc]{}; \draw (0,-1.2) node[ocirc]{}; \draw (9.2,3.5) node[ocirc]{}; \draw (9.2,-1.2) node[ocirc]{};
  \draw[fleche] (-0.35,-1.0) -- node[left]{$v_i$} (-0.35,3.3);
  \draw[fleche] (9.55,-1.0) -- node[right]{$v_o$} (9.55,3.3);
  \draw (0,3.5) -- (1.9,3.5);  \draw[fleche] (0.4,3.7) -- node[above]{$i_i = i_e$} (1.3,3.7);
  \draw (1.9,2.0) -- (1.2,2.0) -- (1.2,0.2) -- (1.9,0.2);
  \draw (0,-1.2) -- (1.9,-1.2);
  \draw (7.4,3.5) -- (9.2,3.5);  \draw[fleche] (8.9,3.7) -- node[above]{$i_o$} (8.0,3.7);
  \draw (7.4,2.0) -- (8.1,2.0) -- (8.1,0.2) -- (7.4,0.2);
  \draw (7.4,-1.2) -- (9.2,-1.2);
\end{circuitikz}
\end{document}
